Este capítulo describe cómo se evalúa JobAI. Primero se enmarca el trabajo en el paradigma de Design Science Research y se reconstruyen los ciclos de diseño-evaluación por los que ya pasó el artefacto. Después se formulan las preguntas de evaluación y se detallan el sistema bajo prueba, el corpus, el conjunto de consultas, los sistemas comparados, el procedimiento de juicios de relevancia, las métricas, el diseño del benchmark de almacenes vectoriales y el análisis estadístico. El capítulo cierra con las amenazas a la validez y las medidas de reproducibilidad.

\section{Enfoque metodológico: Design Science Research}
\label{sec:met-dsr}

Este trabajo sigue el paradigma de Design Science Research (DSR) en la formulación de \citet{Hevner2004}, operacionalizada mediante el proceso de seis actividades de \citet{Peffers2007}. El artefacto es una instanciación de software, JobAI, que ya existía antes de iniciar esta investigación. Por ello se adopta el punto de entrada \emph{centrado en el diseño y desarrollo} que \citet{Peffers2007} prevén para la investigación que parte de un artefacto existente y busca evaluarlo y mejorarlo. La Tabla~\ref{tab:met-peffers} muestra la correspondencia entre las actividades del proceso y el contenido de este documento.

\begin{table}[htbp!]
    \centering
    \caption{Correspondencia entre las actividades de DSR \citep{Peffers2007} y este trabajo}
    \label{tab:met-peffers}
    \begin{tabular}{|p{4.2cm}|p{10cm}|}
        \hline
        \textbf{Actividad} & \textbf{Correspondencia en este trabajo} \\
        \hline
        Identificación del problema y motivación & Limitaciones del emparejamiento léxico en la búsqueda de empleo y ausencia de una evaluación cuantitativa de JobAI (Capítulo~\ref{cap:problema}). \\
        \hline
        Definición de los objetivos de la solución & Objetivo general y objetivos específicos (Capítulo~\ref{cap:problema}), traducidos en preguntas de evaluación medibles (Sección~\ref{sec:met-preguntas}). \\
        \hline
        Diseño y desarrollo & Arquitectura del pipeline RAG y su iteración actual: corrección de la consulta vectorial, migración a Qdrant y corrección de defectos detectados (Capítulo~\ref{cap:artefacto}). \\
        \hline
        Demostración & Ejecución del pipeline sobre un corpus real congelado y un conjunto de consultas de prueba (Secciones~\ref{sec:met-corpus} y~\ref{sec:met-consultas}). \\
        \hline
        Evaluación & Comparación contra búsqueda léxica, medición de fidelidad y latencia, y benchmark pgvector frente a Qdrant (Capítulo~\ref{cap:resultados}). \\
        \hline
        Comunicación & El presente documento. \\
        \hline
    \end{tabular}
\end{table}

\subsection{Ciclos de diseño-evaluación previos}
\label{sec:met-ciclos-previos}

El historial de versiones de JobAI documenta tres iteraciones de diseño-evaluación anteriores a este trabajo, realizadas en abril de 2026 sobre un corpus distinto al actual: unas 3.200 ofertas, mayoritariamente del agregador Adzuna, que dejó de usarse después. Cada iteración se evaluó con 15 consultas y un modelo de lenguaje como juez de relevancia. La Tabla~\ref{tab:met-iteraciones} resume los resultados registrados.

\begin{table}[htbp!]
    \centering
    \caption{Iteraciones previas de diseño-evaluación registradas en el historial de JobAI (15 consultas, $k=5$)}
    \label{tab:met-iteraciones}
    \begin{tabular}{|p{6.2cm}|p{2.6cm}|c|c|c|}
        \hline
        \textbf{Configuración} & \textbf{Juez} & \textbf{P@5} & \textbf{MRR} & \textbf{nDCG} \\
        \hline
        Fragmentación por ventana deslizante & Llama 3.1 8B & 66,7\,\% & 0,656 & 0,702 \\
        \hline
        Fragmentación estructurada (4 tipos) & Llama 3.1 8B & 92,0\,\% & 0,967 & 0,967 \\
        \hline
        Estructurada + ruta de agregación & Llama 3.1 8B & 90,7\,\% & 1,000 & 0,995 \\
        \hline
        Estructurada + ruta de agregación & Modelo comercial & 98,7\,\% & 1,000 & 0,994 \\
        \hline
    \end{tabular}
\end{table}

Estas iteraciones justificaron decisiones de diseño concretas, como la fragmentación estructurada que el Capítulo~\ref{cap:artefacto} describe. Su valor como evidencia, sin embargo, es limitado por cuatro razones:
\begin{itemize}
    \item El conjunto de consultas es pequeño y se diseñó a la vista del corpus.
    \item No existe un sistema de referencia contra el cual comparar.
    \item El juez de relevancia nunca se validó contra juicios humanos.
    \item El corpus ya no corresponde al que la plataforma sirve hoy.
\end{itemize}
La última fila de la Tabla~\ref{tab:met-iteraciones} ilustra la tercera limitación: la misma configuración obtiene una precisión de 90,7\,\% o de 98,7\,\% según el modelo que actúe como juez. La evaluación diseñada en este capítulo responde directamente a estas cuatro limitaciones y constituye, en términos de DSR, la cuarta iteración del ciclo.

\section{Preguntas de evaluación}
\label{sec:met-preguntas}

Los objetivos específicos tercero y cuarto se operacionalizan en cuatro preguntas de evaluación (PE):

\begin{description}
    \item[PE1 (calidad de recuperación).] ¿La recuperación densa de JobAI obtiene mayor precisión y mejor ordenamiento de resultados que la búsqueda léxica, y en qué tipos de consulta se concentra la diferencia?
    \item[PE2 (fidelidad de la generación).] ¿En qué medida las respuestas que genera el pipeline se sostienen en las ofertas recuperadas y responden a la pregunta formulada?
    \item[PE3 (latencia).] ¿Cuánto tarda cada etapa del pipeline (\textit{embedding} de la consulta, recuperación, construcción del contexto y generación) y cuál es la latencia percibida por el usuario?
    \item[PE4 (almacenamiento vectorial).] ¿Cómo se comparan pgvector y Qdrant en la relación entre \textit{recall} y latencia, con y sin filtros estructurados y a distintas escalas de datos?
\end{description}

\section{Sistema bajo evaluación}
\label{sec:met-sistema}

Se evalúa una configuración congelada de JobAI, identificada por el \textit{commit} del repositorio que acompaña a cada corrida de evaluación. La Tabla~\ref{tab:met-config} resume sus parámetros.

\begin{table}[htbp!]
    \centering
    \caption{Configuración del sistema bajo evaluación}
    \label{tab:met-config}
    \begin{tabular}{|p{5cm}|p{9.2cm}|}
        \hline
        \textbf{Componente} & \textbf{Configuración} \\
        \hline
        Modelo de \textit{embeddings} & e5-base-v2 \citep{Wang2022E5}, 768 dimensiones, ejecutado localmente; prefijos \texttt{query:} y \texttt{passage:} \\
        \hline
        Fragmentación & Estructurada (identidad, requisitos, responsabilidades, beneficios y salario); ventana deslizante de 800 caracteres con solapamiento de 200 cuando el análisis estructurado falla \\
        \hline
        Recuperación & 15 fragmentos más similares por coseno, agrupados por oferta; se conservan las 5 ofertas con mayor similitud \\
        \hline
        Índice vectorial & HNSW con $m=16$ y $ef_{construction}=128$ en pgvector y en Qdrant; $ef_{search}=40$ \\
        \hline
        Modelo generador & \texttt{qwen/qwen3.8-27b} servido por Groq, temperatura 0 durante la evaluación \\
        \hline
        Clasificación de consultas & Patrones \textit{regex} más un clasificador basado en LLM restringido a un catálogo cerrado de plantillas SQL \\
        \hline
    \end{tabular}
\end{table}

En producción la temperatura de generación es 0,2. Para la evaluación se fija en 0 y todas las respuestas se almacenan en caché, de modo que una nueva ejecución del análisis no exige regenerarlas. Las consultas se envían sin usuario autenticado; de lo contrario, la personalización por currículo alteraría la consulta y confundiría el efecto de la recuperación con el del perfil del usuario.

\section{Corpus de evaluación}
\label{sec:met-corpus}

El corpus se obtuvo el 10 de octubre de 2026 mediante una ejecución completa del proceso de ingesta de JobAI sobre sus tres fuentes públicas: Remotive, Jobicy y Findwork. El proceso se ejecutó sobre una base de datos dedicada y con el programador de ingesta periódica desactivado, para que el catálogo no cambiara durante la evaluación. El resultado se congeló en una copia de seguridad (\textit{snapshot}) identificada por su resumen SHA-256. Todas las corridas de evaluación parten de esa copia.

El corpus contiene 1.068 ofertas: 895 de Findwork, 162 de Jobicy y 11 de Remotive. De ellas, 682 están marcadas como remotas y 504 informan salario. Las ofertas se dividen en 4.013 fragmentos, todos con su \textit{embedding}, con una longitud media de 467 caracteres. El análisis estructurado con el modelo local Llama 3.1 8B falló por completo en 25 ofertas (2,3\,\%), a las que se aplicó la fragmentación por ventana deslizante, y fue parcial, sin resumen, en otras 231.

El corpus tiene tres sesgos conocidos, que se tienen en cuenta al interpretar los resultados:
\begin{itemize}
    \item Los conectores de Remotive y Jobicy descartan las ofertas sin salario explícito, lo que explica su baja participación.
    \item La API de Findwork interrumpió la descarga al alcanzar su límite de solicitudes, y aun así aporta el 84\,\% del corpus, de modo que este está dominado por una sola fuente.
    \item El normalizador de salarios del artefacto no convierte monedas ni algunos montos por hora, por lo que unas pocas ofertas tienen salarios atípicos. La mediana de los salarios informados es de 175.000 dólares anuales.
\end{itemize}

\section{Conjunto de consultas de prueba}
\label{sec:met-consultas}

El conjunto de prueba contiene 75 consultas en inglés, que es el idioma de las ofertas, distribuidas en cinco categorías de 15 consultas cada una:

\begin{description}
    \item[Exactas.] Consultas cuyos términos aparecen literalmente en las ofertas relevantes, por ejemplo \textit{"senior react developer"}. Es el caso favorable a la búsqueda léxica.
    \item[Desajuste de vocabulario.] Consultas que describen la necesidad con términos distintos a los de las ofertas: sinónimos, títulos equivalentes o descripciones funcionales. Es el caso que motiva este trabajo (Sección~\ref{sec:problema}).
    \item[Con restricciones.] Consultas que combinan un perfil con condiciones sobre atributos estructurados, como salario, modalidad o tipo de contrato.
    \item[Ruidosas.] Consultas con errores ortográficos, abreviaturas o formulación vaga.
    \item[De agregación.] Preguntas estadísticas sobre el catálogo, por ejemplo cuántas ofertas remotas de Python existen. Se evalúan aparte (Sección~\ref{sec:met-metricas}).
\end{description}

Cada consulta se asocia a una \emph{necesidad de información}, es decir, una descripción breve de lo que el usuario busca. A partir de la distribución real de títulos y habilidades del corpus se redactan además 25 consultas en español, traducciones paralelas de consultas en inglés que comparten su necesidad de información. Como la relevancia de una oferta depende de la necesidad y no del idioma en que se formula, las consultas en español reutilizan los juicios de relevancia de su contraparte en inglés. Este subconjunto mide la recuperación entre idiomas, con una salvedad que se anticipa: e5-base-v2 es un modelo entrenado en inglés, por lo que se espera un desempeño inferior en español.

\section{Sistemas comparados}
\label{sec:met-sistemas}

Se comparan cinco sistemas de recuperación sobre el mismo corpus. Todos producen una lista ordenada de 10 ofertas por consulta:

\begin{description}
    \item[ILIKE.] La búsqueda por palabra clave que la plataforma ofrece hoy en su panel de filtros: coincidencia de subcadena, sin distinguir mayúsculas, sobre el título y la empresa. Representa la línea base \emph{de la propia plataforma}.
    \item[BM25.] El modelo probabilístico de relevancia Okapi BM25 \citep{Robertson2009}, con $k_1 = 1{,}2$ y $b = 0{,}75$, sobre el título, la descripción y las habilidades de cada oferta. Representa la línea base léxica \emph{fuerte} de la literatura de recuperación de información. Su inclusión evita que la comparación se haga únicamente contra un sistema débil.
    \item[Denso, pgvector exacto.] Búsqueda exhaustiva por similitud coseno en PostgreSQL. Sirve de referencia para medir el efecto de la búsqueda aproximada.
    \item[Denso, pgvector HNSW.] La misma búsqueda servida por el índice HNSW de pgvector.
    \item[Denso, Qdrant HNSW.] La misma búsqueda servida por Qdrant con parámetros de índice idénticos.
\end{description}

En los tres sistemas densos se recuperan los 50 fragmentos más similares a la consulta y se agregan por oferta, tomando como puntuación de cada oferta la máxima similitud entre sus fragmentos. Es la misma regla de agregación que usa el pipeline en producción. Así la unidad evaluada es la oferta, que es lo que el usuario ve, y no el fragmento.

\section{Juicios de relevancia}
\label{sec:met-juicios}

\subsection{Agrupamiento de resultados}

Los juicios se construyen mediante agrupamiento (\textit{pooling}) de profundidad 10, el procedimiento estándar de las colecciones de prueba en recuperación de información \citep{Sanderson2010}. Para cada consulta se juzga la unión de las 10 primeras ofertas devueltas por cada sistema. Las ofertas que ningún sistema recupera en esas posiciones se consideran no relevantes, una aproximación cuyo efecto se discute en la Sección~\ref{sec:met-validez}.

\subsection{Modelo de lenguaje como juez}

Cada par consulta-oferta del agrupamiento recibe una etiqueta graduada en la escala de cuatro niveles de las evaluaciones TREC, asignada por un modelo de lenguaje:
\begin{itemize}
    \item 0: no relevante.
    \item 1: relacionada.
    \item 2: altamente relevante.
    \item 3: perfectamente relevante.
\end{itemize}
El prompt sigue el diseño de UMBRELA \citep{Upadhyay2024}, una reproducción abierta del procedimiento con el que \citet{Thomas2024} mostraron que un LLM puede predecir las preferencias de evaluadores humanos con una precisión comparable a la de estos. El juez recibe la necesidad de información y una ficha estructurada de la oferta: título, empresa, ubicación, modalidad, salario, habilidades, resumen y un extracto de la descripción.

El uso de un LLM como juez tiene riesgos documentados, entre ellos sesgos sistemáticos, sensibilidad al prompt y posible circularidad cuando el juez y el sistema evaluado comparten modelo \citep{Faggioli2023,Zheng2023}. Para mitigarlos se toman tres medidas:
\begin{itemize}
    \item El juez es un modelo distinto del generador de JobAI.
    \item Durante el desarrollo del arnés de evaluación se usa un modelo local (Llama 3.1 8B), y los resultados definitivos se obtienen con un modelo comercial de mayor capacidad (\pendiente{modelo del juez final}).
    \item Las etiquetas del juez se validan contra juicios humanos, como se describe a continuación.
\end{itemize}

\subsection{Validación humana del juez}

Del conjunto de pares juzgados se extrae una muestra aleatoria del 20\,\%, estratificada por categoría de consulta y por etiqueta del juez, con semilla fija. El autor la etiqueta a ciegas, sin conocer la etiqueta del modelo. La muestra se usa en dos partes:
\begin{itemize}
    \item Una cuarta parte (5\,\% del total) se usa para calibrar el prompt del juez.
    \item Las tres cuartas partes restantes (15\,\% del total) se reservan para medir el acuerdo final, sin volver a modificar el prompt.
\end{itemize}

El acuerdo se mide con el coeficiente kappa de \citet{Cohen1960}, en dos variantes: sobre la relevancia binaria (etiqueta $\geq 2$) y ponderado cuadráticamente sobre la escala completa. Se reporta con un intervalo de confianza bootstrap. Siguiendo la escala de interpretación de \citet{Landis1977}, se considera aceptable un acuerdo al menos sustancial ($\kappa > 0{,}60$).

\section{Métricas}
\label{sec:met-metricas}

\subsection{Calidad de recuperación (PE1)}

Sobre las listas de 10 ofertas se calculan cinco métricas:
\begin{itemize}
    \item La precisión en las posiciones 5 y 10 (P@5, P@10).
    \item La ganancia acumulada descontada normalizada en la posición 10 (nDCG@10) con las etiquetas graduadas \citep{Jarvelin2002}.
    \item El rango recíproco medio (MRR).
    \item El \textit{recall} en la posición 10, relativo al conjunto de ofertas relevantes identificadas en el agrupamiento.
\end{itemize}
Las métricas binarias consideran relevante una oferta con etiqueta $\geq 2$. El cálculo se realiza con la biblioteca \texttt{ir\_measures} \citep{MacAvaney2022}, que unifica las implementaciones de referencia de la comunidad. Los resultados se desagregan por categoría de consulta y por idioma.

Las consultas de agregación no producen una lista de ofertas. Para ellas se miden dos cosas: la exactitud del enrutamiento, es decir, si el clasificador elige la ruta y la plantilla SQL esperadas, y la exactitud de la respuesta frente al resultado de una consulta SQL de referencia escrita a mano.

\subsection{Fidelidad de la generación (PE2)}

Para las consultas de recuperación se generan respuestas mediante la ruta real del pipeline y se calculan dos métricas de RAGAS \citep{Es2024}, empleando la implementación de referencia de los autores:

\begin{itemize}
    \item \textbf{Fidelidad} (\textit{faithfulness}). Un LLM descompone la respuesta en afirmaciones atómicas y verifica cada una contra el contexto recuperado:
    \[
    F = \frac{|V|}{|S|},
    \]
    donde $S$ es el conjunto de afirmaciones de la respuesta y $V \subseteq S$ el de las afirmaciones sostenidas por el contexto.
    \item \textbf{Relevancia de la respuesta} (\textit{answer relevancy}). Un LLM genera $n$ preguntas a partir de la respuesta, y la métrica es la similitud coseno media entre sus \textit{embeddings} y el de la pregunta original:
    \[
    AR = \frac{1}{n}\sum_{i=1}^{n} \cos(\mathbf{e}_q, \mathbf{e}_{q_i}).
    \]
\end{itemize}

\subsection{Latencia (PE3)}

El pipeline se instrumenta para registrar la duración de cada etapa: \textit{embedding} de la consulta, búsqueda vectorial, lectura de las ofertas y construcción del contexto, tiempo hasta el primer \textit{token} generado (TTFT) y tiempo total. Las etapas locales se miden con 30 repeticiones por consulta y por almacén vectorial, tras 5 de calentamiento. La generación, que depende de un servicio externo, se mide con 3 repeticiones sobre un subconjunto de 30 consultas. Se reportan la mediana (p50) y el percentil 95 (p95). La latencia del servicio externo se reporta por separado de la de los componentes propios del artefacto.

\section{Benchmark de almacenes vectoriales (PE4)}
\label{sec:met-benchmark}

La comparación entre pgvector y Qdrant adapta la metodología de ANN-Benchmarks \citep{Aumuller2020}: para cada configuración del parámetro de búsqueda $ef_{search} \in \{16, 32, 64, 128, 256, 512\}$ se mide el \textit{recall}@10 frente a la búsqueda exacta y la latencia de las consultas. El resultado se presenta como curvas de compromiso entre \textit{recall} y latencia, en lugar de un único punto de operación.

\subsection{Datos}

El corpus real es pequeño: unos pocos miles de fragmentos. A esa escala es esperable que ambos motores rindan de forma casi idéntica. Para estudiar el comportamiento a mayor escala se generan conjuntos sintéticos de 100.000 y 500.000 vectores, y de un millón si los recursos de cómputo lo permiten. Cada vector sintético se obtiene interpolando dos vectores reales del corpus y agregando ruido gaussiano escalado por la desviación estándar de cada dimensión, para luego normalizarlo. La magnitud del ruido se calibra de modo que la distribución de similitudes al vecino más cercano se parezca a la del corpus real; la diferencia se cuantifica con el estadístico de Kolmogórov-Smirnov. Cada vector sintético hereda los atributos estructurados de la oferta de la que proviene, de modo que los filtros conservan su correlación con el contenido.

\subsection{Protocolo}

Se emplean unas 1.000 consultas: las del conjunto de prueba más títulos de ofertas del corpus. La verdad de referencia son los 100 vecinos más cercanos exactos, calculados por fuerza bruta, también bajo cada filtro. Para reproducir el patrón de búsqueda mixta vector-relacional estudiado por \citet{Sanca2024}, cada configuración se mide sin filtro y con filtros estructurados de selectividad 1\,\%, 10\,\% y 50\,\%. En pgvector se comparan además los barridos iterativos del índice (\texttt{hnsw.\allowbreak iterative\_scan}), disponibles desde la versión 0.8, frente a su ausencia.

Ambos motores usan el mismo índice HNSW ($m=16$, $ef_{construction}=128$) y se ejecutan uno a la vez, en contenedores con los mismos límites de CPU y memoria. Para cada configuración se registran siete medidas:
\begin{itemize}
    \item \textit{recall}@10;
    \item latencia p50, p95 y p99;
    \item consultas por segundo con 1 y con 8 clientes concurrentes;
    \item completitud, es decir, si se devolvieron $k$ resultados;
    \item tiempo de construcción del índice;
    \item tamaño en disco;
    \item memoria pico.
\end{itemize}
Cada medición se repite tres veces y se reporta la mediana.

\section{Análisis estadístico}
\label{sec:met-estadistica}

Las comparaciones entre sistemas de recuperación son pareadas: todos los sistemas responden las mismas consultas. Siguiendo la comparación de pruebas de significancia en recuperación de información de \citet{Smucker2007}, que desaconseja la prueba de signos y la de Wilcoxon en este contexto, la diferencia en cada métrica se contrasta con la prueba de aleatorización pareada, con 10.000 permutaciones. Se acompaña de un intervalo de confianza bootstrap del 95\,\% para la diferencia de medias. Como se realizan varias comparaciones simultáneas (cada sistema denso frente a cada línea base, por métrica), los valores $p$ se ajustan con el procedimiento secuencial de \citet{Holm1979}. El nivel de significancia es $\alpha = 0{,}05$.

\section{Amenazas a la validez}
\label{sec:met-validez}

\begin{description}
    \item[Validez interna.] Hay tres amenazas principales.
    \begin{itemize}
        \item \emph{El juez automático puede tener sesgos.} Se mitiga con la validación humana y el reporte del kappa.
        \item \emph{El agrupamiento puede favorecer a los sistemas que lo originaron.} Si se añade un sistema después, solo se juzgan sus resultados no evaluados previamente.
        \item \emph{El autor redacta las consultas.} Se mitiga construyéndolas a partir de la distribución real del corpus y por categorías definidas de antemano.
    \end{itemize}
    \item[Validez externa.] Los resultados corresponden a un dominio (empleo remoto, mayoritariamente tecnológico), a un idioma de las ofertas (inglés) y a un único modelo de \textit{embeddings}. No se generalizan a otros dominios ni modelos. El tamaño del corpus real limita las conclusiones sobre escalabilidad; el benchmark sintético las complementa, pero no las sustituye.
    \item[Medición de latencia.] La latencia se mide desde el cliente, de modo que incluye el protocolo de comunicación y la biblioteca cliente de cada motor (\textit{wire protocol} de PostgreSQL frente a gRPC en Qdrant). Es la latencia que percibe una aplicación, pero a escalas pequeñas puede estar dominada por esa sobrecarga y no por el motor. Por eso los resultados se interpretan junto con la escala del conjunto de datos.
    \item[Validez de constructo.] La relevancia juzgada, ya sea por un modelo o por el autor, aproxima pero no equivale a la utilidad que un usuario real percibiría. La evaluación con usuarios queda fuera del alcance (Sección~\ref{sec:problema}).
    \item[Confiabilidad.] Los modelos de lenguaje no son deterministas, incluso a temperatura 0. Además, los proveedores retiran modelos sin previo aviso: durante este trabajo dejaron de estar disponibles los dos modelos generadores configurados originalmente en JobAI (\texttt{qwen3-32b} y \texttt{qwen3.6-27b}), y se adoptó su sucesor, \texttt{qwen3.8-27b}. Ambos riesgos se mitigan registrando cada respuesta y cada veredicto con el identificador exacto del modelo y la fecha, y ejecutando el análisis sobre esa caché.
\end{description}

\section{Reproducibilidad}
\label{sec:met-reproducibilidad}

Cada corrida de evaluación se ejecuta mediante una interfaz de línea de comandos del propio artefacto y escribe sus resultados en un directorio propio con cuatro elementos:
\begin{itemize}
    \item el \textit{commit} del código;
    \item los modelos y parámetros empleados;
    \item el resumen SHA-256 del corpus;
    \item los resultados por consulta.
\end{itemize}
Las versiones de PostgreSQL y pgvector (0.8.2), de Qdrant (1.19.1) y de las bibliotecas de software están fijadas. Los veredictos del juez se almacenan en una caché indexada por modelo, versión del prompt, necesidad de información y oferta, de modo que el análisis puede repetirse sin nuevas llamadas a modelos de lenguaje. El código, el conjunto de consultas, los juicios de relevancia y los resultados se publican junto con este documento. El corpus no se publica, para respetar los términos de uso de las fuentes, pero el procedimiento para reconstruirlo está documentado.
